本文将尝试给出一个Zakharov–Shabat构造的具体例子，
并验证其确实构造了一个可积系统。
关于Zakharov–Shabat构造可以参考\cite{1}中对应部分。
\section{例子}
考虑最简单的情况，只在0处有一阶的pole。记$L=L_0+1/\lambda\cdot L_1$.
假设$P=\lambda L^2$,计算可得
$$M=(P(L))_{-}=1/\lambda\cdot L_1^2$$
$$\frac{dL}{dt}=[L,M]=1/\lambda[L_0,L_1^2]$$

即
\begin{align*}
\dot{L_0}&=0\\
\dot{L_1}&=[L_0,L_1^2]
\end{align*}
因此$L_0$为常量，我们将它对角化：
$$L_0=\begin{pmatrix}u& 0\\0& v\end{pmatrix}$$
设
$$L_1=\begin{pmatrix}a& b\\c& d\end{pmatrix}$$
继续展开
\begin{align*}
\dot{a}&=0\\
\dot{d}&=0\\
\dot{b}&=(a+d)(u-v)b\\
\dot{c}&=(a+d)(v-u)c
\end{align*}

最后变成一个等价于
\begin{align*}
\dot{x}&=x\\
\dot{y}&=-y
\end{align*}
的微分方程。对应的辛形式是$dx\wedge dy$,哈密顿量$xy$。
注意这和一般的Liouville可积系统不一样，后者通常假设哈密尔顿量在任意方向都是增长的（即对任意$\epsilon$, $\{H\le \epsilon\}$紧致）
\begin{thebibliography}
    \bibitem[1]{1}Babelon, O., Bernard, D. \& Talon, M. Introduction to Classical Integrable Systems. (Cambridge University Press,2003)
\end{thebibliography}
